% =====================================================================
% RASC-specific macros
% Shared packages and commands live in ../macros.tex.
%
% BEFORE UNCOMMENTING RASC IN thesis.tex, resolve these conflicts:
%  1. \newcolumntype{C} — ControlInContext defines C with p{} (top-aligned),
%     RASC defines it with m{} (middle-aligned). Rename one (e.g., \Cm or \Cp).
%  2. \sysname — already defined by CoMesh. The \renewcommand below handles it.
%  3. All \newtheorem / \squishlist / \Figure etc. conflicts are already resolved
%     by moving them to ../macros.tex.
% =====================================================================

% Settings for long URLs in references
\def\UrlBreaks{\do\/\do-\do\.}

% --- RASC-only packages ---------------------------------------------
\usepackage{pbalance}   % two-column balance (RASC standalone; harmless in thesis)
\usepackage{comment}    % \begin{comment}...\end{comment} blocks
\usepackage{minted}     % syntax-highlighted code listings
\usepackage{pifont}     % \ding{} symbols (checkmarks, crosses)
\usepackage{siunitx}    % SI units (\SI{}{})
\usepackage{wrapfig}    % text-wrapped figures

% --- Algorithm line number formatting (RASC style) ------------------
\makeatletter
\algrenewcommand\alglinenumber[1]{%
  \makebox[1.5em][r]{\footnotesize #1:}\hspace*{0.25em}%
}
\makeatother

% \newcolumntype{C} is defined in ControlInContext/macros.tex (p{}, top-aligned).
% RASC text files don't use this column type, so no definition is needed here.

% --- RASC check/cross mark symbols ----------------------------------
\newcommand{\cmark}{\ding{51}}
\definecolor{lightgray}{gray}{0.8}
\newcommand{\hxmark}{\textcolor{lightgray}{\ding{51}}}
\definecolor{darkgray}{gray}{0.6}
\newcommand{\hcmark}{\textcolor{darkgray}{\ding{51}}}
\newcommand{\xmark}{\ding{55}}
\definecolor{darkred}{RGB}{153, 0, 0}

\newcommand{\RPCa}{RPC-A}
\newcommand{\RPCs}{RPC-S}
\newcommand{\RPCc}{RPC-C}
\newcommand{\abstraction}{RASC}
\newcommand{\abs}{RASC}
\newcommand{\dynpoller}{our dynamic polling strategy}
\newcommand{\scheduler}{DAG-TL}
\newcommand{\rescheduler}{our rescheduler}
\newcommand{\Rgaragetrash}{R1}
\newcommand{\Rlockdoor}{R1}
\newcommand{\Rfiredetected}{R2}
\newcommand{\Rperimeterlights}{R2}
\newcommand{\Relevatorup}{R3}
\newcommand{\Rpetvacuum}{R3}
\newcommand{\Rprinter}{R4}

% --- RASC theorem style (parenthesized theorem names) ---------------
\newtheoremstyle{thmparen}%
  {3pt}{3pt}%
  {}{}%
  {\bfseries}%
  {.}%
  {0.5em}%
  {\thmname{#1}~\thmnumber{#2}\thmnote{ (#3)}}

\theoremstyle{thmparen}

% Generic "repeat theorem" environment for restating theorems by number
\newtheorem{innercustomgeneric}{\customgenericname}
\providecommand{\customgenericname}{}

\newcommand{\newcustomtheorem}[2]{%
  \newenvironment{#1}[2][]%
  {%
    \renewcommand\customgenericname{#2}%
    \renewcommand\theinnercustomgeneric{##2}%
    \innercustomgeneric[##1]%
  }{\endinnercustomgeneric}%
}

\newcustomtheorem{reptheorem}{Theorem}


% --- \subsubsubsection (RASC uses a 4th heading level) ---------------
% Custom counter that resets with each \subsubsection and numbers as x.y.z.w.
\newcounter{subsubsubsection}[subsubsection]
\renewcommand{\thesubsubsubsection}{\thesubsubsection.\arabic{subsubsubsection}}
\makeatletter
\newcommand{\subsubsubsection}[1]{%
    \refstepcounter{subsubsubsection}%
    \vspace{0.5\baselineskip}%
    \noindent{\normalfont\normalsize\bfseries\thesubsubsubsection\quad#1}\par%
    \vspace{3pt}%
    \@afterindentfalse\@afterheading%
}
\makeatother
\crefname{subsubsubsection}{section}{sections}
\Crefname{subsubsubsection}{Section}{Sections}

% --- Debug utility --------------------------------------------------
\newcommand{\debugbox}[1]{%
  \fbox{\begin{minipage}{\linewidth}#1\end{minipage}}}
